\clearpage\SourceLeaf{314}

\begin{center}
[POTHENOT'S PROBLEM AND THE QUADRILATERAL.]
\end{center}

\begin{center}
[I.]

[CONSTRUCTION FOR POTHENOT'S PROBLEM.]

[1.]
\end{center}

Let three points $a$, $b$, $c$ be represented by the complex
numbers denoted in the same way, and let $\alpha$, $\beta$, $\gamma$ be three
complex numbers such that
\[
\alpha+\beta+\gamma=0.
\]
Set
\[
\alpha A+\beta b+\gamma c=0
\]
\[
\alpha a+\beta B+\gamma c=0
\]
\[
\alpha a+\beta b+\gamma C=0,
\]
and let $A$, $B$, $C$ be the points to which the complex numbers
$A$, $B$, $C$ correspond. Then $aA$,
$bB$, $cC$ intersect at one and the same point $P$, and this is the solution
of \textsc{Pothenot}'s problem.

If $\mathfrak{A}$, $\mathfrak{B}$, $\mathfrak{C}$ are the
(unoriented) azimuths observed at $P$, then
\[
\alpha e^{i\mathfrak{A}}, \qquad \beta e^{i\mathfrak{B}}, \qquad \gamma e^{i\mathfrak{C}}.
\]
are to one another as real numbers.

\Needspace{8\baselineskip}
\begin{center}
REMARK.
\end{center}

\textsc{Gauss} entered this note on the inside of the back cover
of his copy of the work: \textsc{Carnot}, \emph{Geometrie der
Stellung}. Translated by \textsc{Schumacher}. Vol.\ II.\ Altona 1810.

\par\hfill\textsc{Stäckel}.\par
